% page 582 du fac-simile (image p-581 du scan)
% bloc 154.9 x 237.7 mm ; marge gauche 8.0 mm
\pgc{-12.700mm}{0.000mm}{
\l{}
\l{\cel{1}574}
\l{}
\l{\sou{tempestar}. (\sou{netrans.}) Furiar (\sou{aludante la elementi atmosfe}-}
\l{\cel{3}\sou{rala quaz}e\cel{1}\sou{deskatenizita : uragano, fulmino, burasko},}
\l{\cel{3}e c.). -\cel{1}\sou{anke metafore}.- EFIS.}
\l{}
\l{\sou{templo}. Edifico quan onu konsakris a la kulto publika di}
\l{\cel{3}la deajo. - DEFIS.}
\l{}
\l{\sou{temporo}. (\sou{anat.})\cel{2}Regiono laterala di la kapo, inter la}
\l{\cel{3}angulo di la okulo e la parto supra di la orelo.- EFIS}
\l{}
\l{\sou{temporisar}. (\sou{netrans.}) Ajornar, varte di tempo-punto opor-}
\l{\cel{3}tuna. - DEFIS.}
\l{}
\l{\sou{temprar}. (\sou{trans}.) Indutar ye farbo triturita en aquo e mi-}
\l{\cel{3}dilutita en gluo liquida, qua ne esas tam solida kam oleo-}
\l{\cel{3}farbo. - DEFIS.}
\l{}
\l{\sou{tenar}. (\sou{trans}.)Havar en sua manui, inter sua brakii, e\cel{2}c.,}
\l{\cel{3}tale ke (lu) ne povas eskapar, forirar, falar, e c. -}
\l{\cel{3}FIS.}
\l{}
\l{\sou{tenaca}. I. Qua adheras forte ad ulo. - II. Di qua la parti}
\l{\cel{3}inter-adheras forte. - EFIS.}
\l{}
\l{\sou{tenalio}. Utensilo fera, ek du branchi morsoza qui beas e,}
\l{\cel{3}pose, interklemas por sizar e tenar forte. - FIS.}
\l{}
\l{\sou{tenaro}. (\sou{anat.}) Saliajo ek la muskuli, ye la manuo-parto}
\l{\cel{3}avana, extera, supra, apud la polexo. - EFIS.}
\l{}
\l{\sou{tencho}. (\sou{zool.}) Fisho qua vivas en la aquo sen-sala, zoolo-}
\l{\cel{3}gie vicina a karpo. - L.\cel{1}\sou{tinca}. - EFIS.}
\l{}
\l{\sou{tendo}. I. Paviliono ek stofo tensita, quan onu erektas en}
\l{\cel{3}aero ne-enklozita, uzata kom shirmivo (precipue por kampo}
\l{\cel{3}armeala). - II. Kovrilo ek ledro, uzata sur kabrioleto}
\l{\cel{3}od irgaltra veturo sen-tekta, automobila o kaval-tirata.}
\l{\cel{3}- EFIS.}
\l{}
\l{\sou{tende}ncar. (\sou{netrans., ad)}. Sentar en su ta eso-maniero na-}
\l{\cel{3}turale efike da qua ento atraktesas a skopo. - DEFIRS.}
\l{}
\l{\sou{tendino}. (\sou{anat.}) En la muskularo, fasko ek fibri, distingita}
\l{\cel{3}de la muskulo per sua fibri qui esas ne-kontraktebla e}
\l{\cel{3}sua koloro (brilante-blanka). - EFIS.}
\l{}
\l{\sou{tendro}. (\sou{teknol.}) Veturo ligita a lokomotivo por transportar}
\l{\cel{3}la aquo e la karbono ye qui ica alimentesas.- DEFIRS.}
\l{}
\l{\sou{tenebro}. Obskureso absolute kompleta, profunda. II.(\sou{metaf.})}
\l{\cel{3}Eroro, ne-savo, qua impedas la vido di lo exakta, di la}
\l{\cel{3}fakti. - EFIS.}
}
